\clearpage\SourceLeaf{251}\ArithmeticHead{166}{L}{L. SCHLESINGER, ÜBER GAUSS' FUNKTIONENTHEORETISCHE ARBEITEN.}
\begin{center}\textls[120]{b) Das agM. in der Korrespondenz mit \textsc{Schumacher} (1816) und in\\
der Determinatio attractionis.}\end{center}
\noindent\hspace*{\parindent}Während die vorhin erwähnten Abhandlungen sozusagen spontan, kürzere
oder längere Zeit nach ihrem Entstehen von \textsc{Gauss} der Öffentlichkeit übergeben
worden sind, bedurfte es stets eines gewissen äusseren Anstosses, damit
\textsc{Gauss} sich entschloss, mit einem Spezimen aus dem von ihm aufgespeicherten
reichen Schatze der Theorie der transzendenten Funktionen hervorzutreten.
Ein solcher liegt in dem von uns schon wiederholt angeführten Briefe vor,
den \textsc{Schumacher} am 5. April 1816\footnote{Werke X, 1, S. 247, 248.} an \textsc{Gauss} gerichtet hat. \textsc{Schumacher},
der, nach einem nur wenige Jahre dauernden Aufenthalt in Mannheim, wieder
nach Dänemark, und zwar nach Kopenhagen zurückgekehrt war, hat daselbst
den dortigen Professor der Mathematik \textsc{Degen} kennen gelernt. In dem gedachten
Briefe schreibt nun \textsc{Schumacher}, \textsc{Degen} habe ihm erzählt, dass er sich
schon seit längerer Zeit mit einem agM. beschäftige, das, wie \textsc{Schumacher}
meinte, von dem \textsc{Gauss}schen verschieden ist, wobei nämlich für \(a<b\) gebildet
wird \(a'=\frac12(a+b),\ b'=\sqrt{ab}\), \mbox{.\,.\,.\,.\,.} \textsc{Degen} wüsste nur, dass diese Reihe einen
Limes habe, den man sehr bald findet. Er, \textsc{Schumacher}, habe nun den Zusammenhang
dieses Limes mit dem Ellipsenquadranten gefunden und zwar
durch die Formel
\begingroup\makeatletter\tagsleft@true\let\veqno\leqno\makeatother
\begin{equation*}
\frac12\pi\left\{\frac{a'}{M\left(\frac ba\right)}-
\frac{(a-a')\mathfrak{M}\left(\frac ba\right)}{M\left(\frac ba\right)}\right\}
=\text{Quadrans Ellipseos}\left(\frac{b'}{a'}\right),\tag{7}
\end{equation*}
\endgroup
\noindent wo \(\left(\frac{b'}{a'}\right)\) eine Ellipse bedeutet, deren Achsen \(1\) und \(\frac{b'}{a'}\) sind, \(M\left(\frac ba\right)\) das agM.,
und
\begingroup\makeatletter\tagsleft@true\let\veqno\leqno\makeatother
\begin{equation*}
\mathfrak{M}\left(\frac ba\right)=\frac{a-a'}{2a'}+
\frac{a-a'}{2a'}\mathbin{.}\frac{a'-a''}{2a''}+
\frac{a-a'}{2a'}\mathbin{.}\frac{a'-a''}{2a''}\mathbin{.}\frac{a''-a'''}{2a'''}+\ldots\tag{8}
\end{equation*}
\endgroup
\noindent gesetzt wurde. \textsc{Gauss} antwortet (April 1816, Werke X, 1, S. 247, 248), das
agM., mit dem \textsc{Degen} sich beschäftigt, sei dasselbe wie das, mit dem er sich
»seit 1791 beschäftigt habe und jetzt einen ziemlichen Quartband darüber
schreiben könne«; er erinnert \textsc{Schumacher} daran, dass dieser 1808, als er sich
\par\clearpage\SourceLeaf{252}\ArithmeticHead{167}{R}{VIIIb. KORRESPONDENZ ÜBER DAS AGM.}
\noindent studienhalber in Göttingen aufhielt, die 1800 begonnene Abhandlung über
das agM. in dem Handbuch »Varia« (15, Ba, siehe Werke III, S. 361) gelesen
habe. Er, \textsc{Gauss}, werde in dem zweiten Teil der \textit{Disquisitiones} von 1812,
den er bald zu geben hoffe, »einen Teil seiner Untersuchungen über das agM.
bekannt machen«.

Auffallend ist der gereizte Ton dieses \textsc{Gauss}schen Briefes. \textsc{Schumacher}
hatte sich offenbar das, was er seinerzeit (1808) über die Media von \textsc{Gauss}
gelesen hatte, wohl eingeprägt, aber er wusste nicht mehr, dass er es gelesen
hatte und hielt es für seine eigene Erfindung. Solche Fälle kommen bekanntlich
nicht selten vor. Besonders frappant ist die \textls[120]{vollständige} Übereinstimmung
der von \textsc{Schumacher} gegebenen Formel (7) für den Ellipsenquadranten
mit der im \textit{Leiste} von \textsc{Gauss} aufgezeichneten (Werke X, 1, S. 178,
Gln. (14), (15)\footnote{Diese Gleichungen enthalten einige Schreibfehler; vergl. darum die berichtigten Formen in den Bemerkungen, ebenda S. 265.}). Da es wohl ausgeschlossen ist, dass \textsc{Gauss} diese Formel erst
nach dem Empfang des \textsc{Schumacher}schen Briefes in den \textit{Leiste} eingetragen
hat, so wird man diese Übereinstimmung entweder dem »Zufall« zuschreiben
oder annehmen müssen, dass \textsc{Schumacher} die Formel etwa 1808 bei \textsc{Gauss}
gesehen und sie dann 1816 gleichsam aus dem »Unterbewusstsein« reproduziert
hat. \textsc{Gauss} hat sich augenscheinlich gerade über diese Formel geärgert, denn
er betont, dass \textsc{Schumacher} in dem Handbuche die Formel (Werke III, S. 373)
\begingroup\makeatletter\tagsleft@true\let\veqno\leqno\makeatother
\begin{equation*}
d\operatorname{Med.}(x,y)=\operatorname{Med.}(x,y)\times
\left\{\begin{gathered}
\frac{dx}{x}\left(\frac12+\frac14\frac{x-x'}{x'}+\frac18\frac{x-x'}{x'}\mathbin{.}\frac{x'-x''}{x''}+\mathrm{etc.}\right)\\
\frac{dy}{y}\left(\frac12-\frac14\frac{x-x'}{x'}-\frac18\frac{x-x'}{x'}\mathbin{.}\frac{x'-x''}{x''}+\mathrm{etc.}\right)
\end{gathered}\right\},\tag{9}
\end{equation*}
\endgroup
\noindent in der die von \textsc{Schumacher} mit \(\mathfrak{M}\left(\frac ba\right)\) bezeichnete Reihe (8) auftritt, selbst
»restituiert« habe (vergl. die Fussnote \textsuperscript{117)}).

Allerdings steht in diesem Handbuche explizite nichts vom Ellipsenquadranten,
aber \textsc{Schumacher} schreibt in seiner Antwort (vom 8. Juni 1816,
\textit{Briefwechsel Gauss-Schumacher} I, S. 127): »\mbox{.\,.\,.} wie Ihre elementaren Media,
über die ich ein paar Zettel von Ihrer Hand inter {\fontspec{times.ttf}κειμηλια} bewahre \mbox{.\,.\,.}«\footnote{In den im Gaussarchiv vorhandenen \textit{Gaussiana} \textsc{Schumacher}s finden sich Aufzeichnungen über diese »elementaren Media« teils von \textsc{Gauss}', teils von \textsc{Schumacher}s Hand, vergl. Werke X, 1, S. 234—237. Vergl. zu der ganzen Kontroverse die Bemerkungen Werke X, 1, S. 284—286 und S. 265.};
er hat also seinerzeit (1808) auch noch andere Handschriften von \textsc{Gauss} als
\par\clearpage\SourceLeaf{253}\ArithmeticHead{168}{L}{L. SCHLESINGER, ÜBER GAUSS' FUNKTIONENTHEORETISCHE ARBEITEN.}
\noindent das Handbuch eingesehen. Vielleicht hat auf diesen der Quadrans Ellipseos
gestanden, vielleicht waren es jene verlorenen Uraufzeichnungen (vergl. oben
S. 42), die den Leistenotizen über das agM. zugrunde lagen?

In diesem Antwortschreiben \textsc{Schumacher}s sind auch die literarischen Angaben
von Interesse (a. a. O. S. 127), die \textsc{Schumacher} über die von ihm zu
Rate gezogenen Hilfsmittel macht. Er zitiert nämlich eine Abhandlung von
\textsc{Woodhouse} aus den Philosophical Transactions, 1804 II, in der eine Abhandlung
von \textsc{Wallace} exzerpiert sei, und wo die »bekannte Umformung von
\(\int dx\left(\frac{1-eexx}{1-xx}\right)^{\frac12}\) auf ähnliche Formen, wo \(e\) immer kleiner wird, ins Unendliche
fortgesetzt wird«. Es ist dies also die \textsc{Landen}sche Transformation\footnote{Die gedachte Abhandlung von \textsc{Wallace} und eine bei \textsc{Woodhouse} ebenfalls exzerpierte von \textsc{Ivory} befinden sich in den Bänden V. bezw. IV. der Transactions of the Royal Society of Edinburgh. Sie enthalten in bezug auf die \textsc{Landen}sche Transformation nichts, was wesentlich über \textsc{Lagrange}s Abhandlung von 1784/85 (vergl. oben S. 12 ff.) hinausgeht.}. Ob
\textsc{Gauss} diese Literatur schon früher gekannt, oder ob er sie erst aus diesem
Briefe \textsc{Schumacher}s kennen gelernt hat, muss dahingestellt bleiben. Auf eine
andere Stelle dieses \textsc{Schumacher}schen Briefes, die sich auf \textsc{Degen} bezieht,
kommen wir weiter unten zurück.

Die von \textsc{Gauss} ausgesprochene Absicht, den zweiten Teil der \textit{Disquisitiones}
mit einem Teil seiner Untersuchungen über das agM. bald zu geben, kommt
wieder nicht zur Ausführung. Jedoch scheint er — nach längerer Pause —
bald darauf die Untersuchungen über lemniskatische Funktionen wieder aufgenommen
zu haben, wenigstens stammen die Werke III, S. 421, Art. [14]
abgedruckten, dem Handbuche 19, Be (S. 104) entnommenen Aufzeichnungen
aus der Zeit nach dem 20. Februar 1817 (vergl. die Angabe von \textsc{Schering},
Werke III, S. 494). Sie enthalten die Additionstheoreme der Thetafunktionen
im lemniskatischen Falle und das Additionstheorem für das Integral zweiter
Gattung \(\int(\sin\operatorname{lemn}X)^2dX\). Bemerkenswert ist das erstmalige Auftreten der
Zeichen \(P,Q,R,S\) für die vier Thetafunktionen, hier natürlich nur im lemniskatischen
Falle. Jedenfalls verdanken wir es aber dieser Korrespondenz mit
\textsc{Schumacher}, dass \textsc{Gauss} bei der nächsten Gelegenheit wenigstens die Grundzüge
seiner bis dahin streng geheim gehaltenen Theorie des agM. bekannt
macht, in der Abhandlung \textit{Determinatio attractionis quam in punctum quodvis
positionis datae exerceret planeta, si eius massa per totam orbitam ratione temporis}
\par\clearpage\SourceLeaf{254}\ArithmeticHead{169}{R}{VIIIb. DIE DETERMINATIO ATTRACTIONIS.}
\noindent\textit{quo singulae partes describuntur uniformiter esset dispertita}\footnote{Vorgelegt 17. Januar 1818, Comm. Soc. reg. sc. Gott. rec. vol. IV, Werke III, S. 331, \textit{Anzeige} ebenda, S. 357, deutsche Übersetzung bei \textsc{Geppert}, a. a. O. \textsuperscript{83)}.}. Über die Wichtigkeit
des in dieser Arbeit behandelten astronomischen Problems spricht sich
\textsc{Gauss} in einem Briefe an \textsc{Olbers} vom 11. Januar 1818 aus\footnote{\textit{Briefwechsel Olbers-Gauss} I, S. 679, vergl. \textsc{M. Brendel}, Werke XI, 2, Abh. 3, S. 253, 254.}, wo er auch
angibt, dass er sich »bisher« mit der Arbeit an dieser Abhandlung beschäftigt
habe.

Der auf das agM. bezügliche Teil dieser Abhandlung (Art. 16—19) bringt
zuvörderst den elementaren Algorithmus für zwei positive Zahlen \(m,n\). Der
Beweis, dass \(\mu=M(m,n)\) dem reziproken Wert des vollständigen elliptischen
Integrals erster Gattung
\[
\int_0^{2\pi}\frac{dT}{2\pi\sqrt{m^2\cos^2T+n^2\sin^2T}}
\]
gleich ist, wird durch Anwendung der uns bereits bekannten, sog. \textsc{Gauss}schen
Transformation (26) Abschnitt IV (oben S. 109, vergl. auch S. 131) geführt,
die das Differential
\[
\frac{dT}{\sqrt{m^2\cos^2T+n^2\sin^2T}}
\]
invariant lässt. Bildet man durch Wiederholung dieser Transformation die
Folge\linebreak \(T',T'',T''',\ldots\), so strebt diese einem Grenzwerte \(\theta\) zu, und es ist
\[
\frac\theta\mu=\int_0^T\frac{dT}{\sqrt{m^2\cos^2T+n^2\sin^2T}}.
\]
Auch für ein Integral zweiter Gattung mit \textls[120]{variabler} oberer Grenze gibt
\textsc{Gauss} eine Darstellung (Art. 18, vergl. oben S. 109, 110). In der \textit{Anzeige} fügt
\textsc{Gauss} noch hinzu, dass der auf das vollständige Integral erster Gattung bezügliche
Satz sich auch so aussprechen lässt, dass in der Entwicklung
\[
\frac1{\sqrt{\alpha+\beta\cos\psi}}=A+B\cos\psi+C\cos2\psi+\mathrm{etc.}
\]
\(\frac1A\) gleich \(M(\sqrt{\alpha+\beta},\sqrt{\alpha-\beta})\) ist, was direkt an den oben erwähnten Brief an
\textsc{Schumacher} (Werke X, 1, S. 247) erinnert. Das gleiche gilt von dem Satze,
den \textsc{Gauss} in der \textit{Anzeige} als »ein zweites, eben so wichtiges Theorem« bezeichnet,
\HistoricalFoot{X 2 Abh. 2.}{22}
\par\clearpage\SourceLeaf{255}\ArithmeticHead{170}{L}{L. SCHLESINGER, ÜBER GAUSS' FUNKTIONENTHEORETISCHE ARBEITEN.}
\noindent und wonach das vollständige Integral zweiter Gattung
\[
\begin{aligned}
&\int_0^{2\pi}\frac{\cos2T\,dT}{2\pi\sqrt{m^2\cos^2T+n^2\sin^2T}}\\
&\quad=-\frac1\mu\left[2\frac{m'^2-n'^2}{m^2-n^2}
+4\frac{m''^2-n''^2}{m^2-n^2}
+8\frac{m'''^2-n'''^2}{m^2-n^2}+\ldots\right]
\end{aligned}
\]
ist. In der \textit{Anzeige} wird sogar noch besonders die Darstellung des Ellipsenumfangs
durch die Reihe
\[
\frac{2\pi}\mu\bigl(m'^2-2(m''^2-n''^2)-4(m'''^2-n'''^2)-\ldots\bigr)
\]
angegeben, die im wesentlichen auf \textsc{Schumacher}s Formel (7) hinauskommt,
gleichsam als wollte \textsc{Gauss} hierdurch sein Eigentumsrecht an dieser Formel
ausdrücklich betonen.

Den Ausspruch am Schluss der \textit{Anzeige}, wonach \textsc{Gauss} »diese Resultate,
wie er sie schon vor vielen Jahren unabhängig von ähnlichen Untersuchungen
von \textsc{Lagrange} und \textsc{Legendre} gefunden hat, in ihrer ursprünglichen Form
darstellen zu müssen geglaubt hat«, wird man, was die »ursprüngliche Form«
anlangt, cum grano salis aufzufassen haben, denn wir wissen, dass \textsc{Gauss} jedenfalls
den Zusammenhang des agM. mit dem vollständigen Integral erster Gattung
nicht auf dem in dieser Abhandlung angegebenen Wege gefunden hat.
